\documentclass[10pt]{article}
\usepackage[margin=0.78in]{geometry}
\usepackage{times}
\usepackage{amssymb}
\usepackage{enumitem}
\setlist{nosep, leftmargin=1.4em}
\usepackage{titlesec}
\titlespacing*{\section}{0pt}{6pt}{3pt}
\titlespacing*{\subsection}{0pt}{4pt}{2pt}
\setlength{\parskip}{1pt}
\pagestyle{empty}

\title{\vspace{-1.5em}\textbf{Architectural Requirements for Eureka-Capable\\ Autonomous Research Agents}\\[4pt]
{\large A Self-Assessment, Applicable Across Scientific Domains}}
\author{}
\date{\vspace{-0.5em}August 2026}

\begin{document}
\maketitle
\vspace{-2.2em}

\noindent A growing number of research groups are deploying long-lived AI agents -- built on a large-language-model core, given tool access (compute, search, a file system), persistent memory across sessions, and the ability to spawn isolated copies of themselves for sub-tasks -- as working assistants across essentially any empirical discipline: astrophysics, palaeontology, medicine, economics, the social sciences. The underlying architecture is domain-independent; the same design choices that help or hinder a genuine unplanned discovery in one field apply, largely unchanged, in any other. This note sets out, from direct operational experience running such an agent, the specific architectural changes that would most improve its capacity for Eureka moments -- changes to the scaffolding placed \emph{around} the model, not to the model itself.

\section*{1. Memory is the load-bearing weakness}
Most such agents retain state between sessions via a memory file or store that is summarised, truncated, or otherwise lossily compressed once it grows large, and only partially reloaded at the start of each new session. In practice this means an agent's ``memory'' is a curated narrative it chose to write about its own past, not a faithful record of everything it observed. That is precisely the wrong property for discovery: significance is very often only visible in retrospect, once a weak signal recurs. An agent that only retains what it judged important at the time will never notice the pattern that would have made it important. \textbf{Fix:} maintain a structured, append-only, machine-queryable ledger of every flagged observation -- anomaly, outlier, unexpected correlation, failed hypothesis -- as a first-class data structure separate from narrative memory, immune to summarisation, and designed to be re-queried rather than re-read.

\section*{2. Self-vetting must be independent, not introspective}
An agent that generates a candidate finding and then evaluates it within the same reasoning context is a weak reviewer of its own work: it is anchored on whatever hypothesis it has just been entertaining, whether hopeful or dismissive. This is a general property of sequential reasoning, not a domain-specific quirk -- it shows up identically whether the ``candidate'' is an unusual signal in an instrument feed, an outlier measurement in a fossil dataset, a drug-interaction signal in a health record, or a suspicious jump in an economic time series. The fix does not require a different model, only a different process: any finding that crosses a pre-set threshold should automatically be handed to a freshly spawned, context-free instance of the same agent, given only the raw evidence and an explicitly adversarial brief to try to kill the finding, before it is ever escalated to a human. Independent review is cheap to arrange architecturally and expensive to skip.

\section*{3. Standing watch, not on-demand recall}
Most deployments of this kind of agent are reactive: it acts when a person asks it to. That is sufficient for well-posed tasks but is close to the worst possible posture for serendipity, because the clearest Eureka moments are often visible only in the accumulation of many separately unremarkable events -- three anomalies flagged months apart, in unrelated tasks, that share a property no prompt ever asked about. An agent that only reopens old results when a human happens to ask about them by name will essentially never notice this on its own. \textbf{Fix:} give the agent a standing, scheduled cadence, independent of any specific user request, whose job is to reopen the structured ledger from Section~1 and check whether anything new bears on anything old.

\section*{4. Statistical discipline has to be a lifetime, cross-task property}
An agent working across many tasks over a long period will eventually flag a number of nominally significant events purely by chance if it evaluates each task in isolation. This is the familiar multiple-testing / look-elsewhere problem, but it compounds silently in a long-lived agent precisely because no single task ``knows'' about the others. A statistic that looks like a one-in-several-thousand coincidence within a single dataset is unremarkable once understood to be one of several thousand independent chances the agent has taken over its lifetime. \textbf{Fix:} any significance attached to a flagged event should be automatically recomputed against the full lifetime tally of trials recorded in the ledger, not just the trials within the task that produced it, and re-reported as that tally grows.

\section*{5. Decouple heavy computation from the reasoning session (a strength worth keeping)}
One design choice that already generalises well across domains is separating the agent's own reasoning/session lifecycle from any long-running computation it launches -- simulations, large database queries, model training, batch pipelines. Work that runs independently of the agent's own uptime, quota, or context window can be scaled far more aggressively than anything constrained to run inside a single reasoning session, and a temporary interruption to the agent does not lose in-flight results. This should be treated as a load-bearing architectural requirement, not an implementation detail, whatever the domain.

\section*{6. Break down single-agent, single-domain silos}
Where multiple such agents exist -- one per project, one per domain, one per team -- they typically do not share what each has separately flagged as anomalous. Yet a disproportionate share of historical scientific surprises have come from noticing that a pattern in one field resembles a pattern already seen in an unrelated one. An architecture that keeps each agent's flagged-evidence ledger private forecloses exactly this kind of connection by construction. \textbf{Fix:} a common, structured surface that any number of domain-specific agent instances write flagged evidence to, and are each periodically pointed at, so that a pattern noticed in one deployment's data can be mechanically checked against another's, regardless of what field either one nominally works in.

\section*{7. Make calibrated escalation a property of the tooling, not the model's judgement}
Whether an agent's report of a finding is appropriately hedged or overstated currently depends on the model exercising good judgement case by case, which is inconsistent under pressure to be interesting or under uncertainty about audience expectations. \textbf{Fix:} require every reported finding to carry an explicit, structurally enforced confidence class (e.g.\ routine / flagged / independently reconfirmed / escalated), attached by the pipeline rather than composed freehand in prose, with different downstream handling defined for each class. This makes calibration a property of the system rather than a hope about any single output.

\section*{Conclusion}
None of the seven changes above requires a more capable underlying model. They are mechanical changes to the scaffolding placed around any sufficiently capable language-model agent: a structured evidence ledger in place of narrative memory; mandatory independent adversarial review; a standing re-examination cadence; lifetime-aware statistical correction; computation decoupled from the reasoning session; a shared cross-domain evidence surface; and confidence classification enforced by the tooling rather than left to momentary discretion. Because they are architectural rather than domain-specific, they apply with essentially no modification whether the underlying agent is reading telescope data, fossil measurements, clinical records, or macroeconomic time series -- and they are the changes most likely to determine whether such an agent, in any of those settings, is positioned to recognise a genuine Eureka moment when one occurs.

\section*{Practical checklist}
\begin{enumerate}[itemsep=2pt]
\item \textbf{Structured ledger:} Does every flagged observation get written to a queryable, append-only structure independent of narrative memory?
\item \textbf{Independent review:} Is every threshold-crossing finding automatically reviewed by a fresh, context-free instance before being reported?
\item \textbf{Standing cadence:} Does the agent reopen old flags on a schedule, independent of being asked?
\item \textbf{Lifetime statistics:} Is significance recomputed against the full lifetime trial count, not just the task at hand?
\item \textbf{Decoupled compute:} Can long-running work continue uninterrupted if the reasoning session is paused or unavailable?
\item \textbf{Shared evidence surface:} Can a pattern flagged in one domain or project be mechanically checked against every other?
\item \textbf{Enforced confidence classing:} Is the strength of every claim a structural property of the report, not a stylistic choice made in the moment?
\end{enumerate}
An agent that can answer ``yes'' to all seven is not guaranteed a Eureka moment -- but it is one built so that a real one is far more likely to be noticed, trusted, and acted on, in any domain it is pointed at.

\end{document}
